\addcontentsline{toc}{chapter}{Abstract}\vspace{-1cm}
%Border
\begin{tikzpicture}[remember picture, overlay]
  \draw[line width = 4pt] ($(current page.north west) + (0.75in,-0.75in)$) rectangle ($(current page.south east) + (-0.75in,0.75in)$);
\end{tikzpicture}


%Highlights of significant contributions: One page with 3 to 4 paragraphs\\

Fluoride contamination in drinking water is a major environmental concern due to its harmful health effects. Conventional fluoride removal methods often involve high cost and complex processing, creating the need for low-cost and eco-friendly alternatives. In the present work, coconut fiber was selected as the base adsorbent material and surface modification using Glycidyl Trimethyl Ammonium Chloride (GTMAC) was carried out to improve fluoride adsorption characteristics. 


The main objective of this work was to prepare and functionalize coconut fiber for fluoride adsorption studies and to develop a systematic experimental design using Design of Experiments (DOE). Biomaterial preparation involved drying, grinding, and sieving of coconut fiber, followed by GTMAC functionalization under alkaline conditions. Central Composite Design (CCD) was used for planning adsorption studies by varying parameters such as pH, adsorbent dose, fluoride concentration, and contact time. 


Approximately 30 CCD experimental runs were generated using Design-Expert software for adsorption studies and optimization. Preliminary Langmuir adsorption model diagnostics carried out during Phase 2 showed acceptable prediction performance with an R² value of 0.8158 and RMSE value of 0.5278 mg/g, indicating good agreement between actual and predicted adsorption values. 

 

\pagebreak